% Standalone problem document, generated from the adjacent README.md.
% Compile with XeLaTeX. Mathematical content is maintained in Markdown.
\documentclass[11pt,a4paper]{article}
\usepackage[fontsize=10.5pt]{scrextend}
\usepackage[margin=22mm,headheight=15pt,headsep=9mm,footskip=11mm]{geometry}
\usepackage{fontspec}
\setmainfont{texgyrepagella}[Extension=.otf,UprightFont=*-regular,BoldFont=*-bold,ItalicFont=*-italic,BoldItalicFont=*-bolditalic]
\setsansfont{texgyreheros}[Extension=.otf,UprightFont=*-regular,BoldFont=*-bold,ItalicFont=*-italic,BoldItalicFont=*-bolditalic]
\setmonofont{texgyrecursor}[Extension=.otf,UprightFont=*-regular,BoldFont=*-bold,ItalicFont=*-italic,BoldItalicFont=*-bolditalic,Scale=MatchLowercase]
\usepackage{amsmath,amssymb}
\usepackage{unicode-math}
\setmathfont{texgyrepagella-math.otf}
\usepackage{xcolor,microtype,enumitem,needspace,fancyhdr}
\usepackage{bookmark}
\definecolor{ink}{HTML}{17344B}
\definecolor{accent}{HTML}{197D80}
\definecolor{muted}{HTML}{596773}
\hypersetup{colorlinks=true,linkcolor=accent,urlcolor=accent,pdftitle={MI-08 - Minimum
number of orthogonal conjugations for diagonal
pinching},pdfauthor={Open Problems in Numerical Linear Algebra}}
\urlstyle{same}
\setlength{\parindent}{0pt}
\setlength{\parskip}{5pt plus 1pt minus 1pt}
\linespread{1.04}
\setlength{\emergencystretch}{3em}
\widowpenalty=10000
\clubpenalty=10000
\displaywidowpenalty=10000
\allowdisplaybreaks[1]
\setlist{nosep,leftmargin=1.5em}
\providecommand{\tightlist}{\setlength{\itemsep}{0pt}\setlength{\parskip}{0pt}}
\providecommand{\pandocbounded}[1]{#1}
\pagestyle{fancy}
\fancyhf{}
\fancyhead[L]{\sffamily\footnotesize\color{muted}OPEN PROBLEMS IN NUMERICAL LINEAR ALGEBRA}
\fancyhead[R]{\sffamily\bfseries\footnotesize\color{accent}MI-08}
\fancyfoot[L]{\parbox[b]{0.9\textwidth}{\sffamily\fontsize{8}{10}\selectfont\color{muted}Editorial ratings; a bounded search cannot certify that a problem remains unsolved.\\Literature check: 2026-09-14}}
\fancyfoot[R]{\sffamily\footnotesize\color{muted}\thepage}
\renewcommand{\headrulewidth}{0.3pt}
\renewcommand{\headrule}{\hbox to\headwidth{\color{accent}\leaders\hrule height \headrulewidth\hfill}}
\makeatletter
\renewcommand{\section}{\Needspace{5\baselineskip}\@startsection{section}{1}{0pt}{12pt plus 2pt minus 2pt}{4pt}{\sffamily\bfseries\large\color{ink}}}
\renewcommand{\subsection}{\Needspace{4\baselineskip}\@startsection{subsection}{2}{0pt}{9pt plus 1pt minus 1pt}{3pt}{\sffamily\bfseries\normalsize\color{ink}}}
\makeatother
\setcounter{secnumdepth}{0}
\begin{document}
{\sffamily\small\color{muted}Matrix inequalities and norms\par}
\vspace{3pt}
{\raggedright\hyphenpenalty=10000\exhyphenpenalty=10000\sffamily\bfseries\fontsize{21}{24}\selectfont\color{ink}Minimum
number of orthogonal conjugations for diagonal pinching\par}
\vspace{7pt}
{\sffamily\small\textbf{Difficulty:} extreme\quad\textbf{Importance:} interesting
to the community\par}
{\sffamily\small\bfseries\color{accent}Status: Partially resolved\par}
\vspace{4pt}
{\color{accent}\hrule height 0.8pt}
\vspace{6pt}
\textbf{Rating rationale:} An exact all-dimension optimum imposes rigid
real orthogonality constraints beyond a general upper bound; it would
clarify matrix averaging and majorization decompositions.

\subsection{Partial result ---
2026-09-11}\label{partial-result-2026-09-11}

\textbf{Partial result by Matthew J. Colbrook}, Department of Applied
Mathematics and Theoretical Physics, University of Cambridge.
\textbf{Independent proof review: PASS for the stated partial result.}

The fixed and adaptive orthogonal pinching lengths both equal the least
row count \(h(d)\) of a sign matrix \(H\) with \(H^TH=h(d)I_d\). In
particular, the exact length is 12 for \(9\le d\le12\).

\textbf{Still open:} The general value of \(h(d)\) is undetermined; the
all-dimension optimization remains open and includes Hadamard-order
existence questions. The ratings apply to this surviving question.

\textbf{Primary manuscript:}
\href{https://github.com/ajt60gaibb/OpenProblemsInNLA/blob/main/matrix-inequalities-and-norms/MI-08/solution.pdf}{complete
proof PDF},
\href{https://github.com/ajt60gaibb/OpenProblemsInNLA/blob/main/matrix-inequalities-and-norms/MI-08/solution.tex}{standalone
TeX}, Theorem 1.1 and its proof;
\href{https://github.com/ajt60gaibb/OpenProblemsInNLA/blob/main/matrix-inequalities-and-norms/MI-08/solution.md}{authorship
and scope}. The
\href{https://github.com/ajt60gaibb/OpenProblemsInNLA/blob/main/references/colbrook-matrix-2026-09-11/verification/reviews/MI-08-review.md}{independent
review} checks the full original argument and records its hash. The
draft was AI-assisted; this is independent agent verification, not
external human peer review or formal certification.
\href{https://github.com/ajt60gaibb/OpenProblemsInNLA/blob/main/references/colbrook-matrix-2026-09-11/README.md}{Submission
record}.

\subsection{Lean proof and verification evidence ---
2026-09-14}\label{lean-proof-and-verification-evidence-2026-09-14}

\textbf{The stated fixed-list partial result is Lean verified; the
canonical problem remains Partially resolved.} The fixed-list
feasibility characterization by rectangular sign matrices is proved for
every positive dimension and row count. The rank and divisibility
obstructions and an explicit order-twelve sign matrix establish that the
actual minimum fixed pinching length is 12 for every dimension from 9
through 12. The source's adaptive comparison is not formalized, and the
all-dimension optimum remains open.

\href{https://github.com/sidneyholden1/OpenProblemsInNLA/tree/a66e142cbe7907db6594ef57c7f542c3d33ad704/matrix-inequalities-and-norms/MI-08/lean}{Immutable
formal proof}. The checked exports are
\textit{NLA.\allowbreak{}MI08.\allowbreak{}fixed\_\allowbreak{}sign\_\allowbreak{}equivalence},
\textit{NLA.\allowbreak{}MI08.\allowbreak{}design\_\allowbreak{}obstructions},
\textit{NLA.\allowbreak{}MI08.\allowbreak{}hadamard\_\allowbreak{}twelve}, \textit{NLA.\allowbreak{}MI08.\allowbreak{}finite\_\allowbreak{}minimums}.
Lean 4.33.1 uses pinned Mathlib \textit{0df444a3} and LeanCert
\textit{621a43d7};
\href{https://github.com/sidneyholden1/OpenProblemsInNLA/tree/codex/lean-mi08/matrix-inequalities-and-norms/MI-08/lean/README.md}{complete
pins and reproduction instructions} are retained. All four exported
transitive axiom closures contain only \textit{propext},
\textit{Classical.choice} and \textit{Quot.sound}.

\textbf{Formalization:} Sidney Holden, Center for Computational Biology,
Flatiron Institute, Simons Foundation, with OpenAI Codex assistance.
Matthew J. Colbrook retains mathematical authorship. Two independent
statement reviews preceded implementation, and two independent final
proof reviews passed. Kernel-only LeanCert and
\href{https://github.com/sidneyholden1/OpenProblemsInNLA/actions/runs/34862187226}{fresh
isolated Linux Comparator and kernel verification} passed, including
rejection controls.
\href{https://github.com/sidneyholden1/OpenProblemsInNLA/tree/codex/lean-mi08/matrix-inequalities-and-norms/MI-08/lean/verification/linux-2026-09-14/OPERATIONAL-REVIEW.md}{Original
artifact, exact checked inputs and independent operational audit} are
retained. This is AI-assisted formal verification, not external human
peer review or source-author endorsement.

\newpage

\subsection{Problem statement}\label{problem-statement}

For each integer \(d\ge1\), define
\(\Delta(X)=\mathop{\mathrm{diag}}\nolimits(x_{11},\ldots,x_{dd})\) for
\(X=(x_{ij})\in\mathbb R^{d\times d}\). Determine the integer

\[\varphi(d)=\min\left\{q\ge1:\ \exists U_1,\ldots,U_q\in O(d)\ \forall X\in\mathbb R^{d\times d},\quad
\Delta(X)=\frac1q\sum_{i=1}^qU_iXU_i^T\right\},\]

where \(O(d)=\{U\in\mathbb R^{d\times d}:U^TU=I_d\}\). The same list of
orthogonal matrices must work for every \(X\). The minimum is finite:
\(q=2^{\lceil\log_2d\rceil}\) is always possible.

\subsection{Why it matters}\label{why-it-matters}

Diagonal extraction is a basic matrix operation. Expressing it by the
shortest average of orthogonal similarities gives an exact decomposition
complexity and sharpens real versions of majorization-based norm
estimates.

\subsection{References}\label{references}

\begin{enumerate}
\def\labelenumi{\arabic{enumi}.}
\tightlist
\item
  J.-C. Bourin and E.-Y. Lee, \emph{Averages over matrix unitary orbits
  and spectral order}, arXiv:2606.15624v2 (18 June 2026), Lemma 4.1 and
  Question 4.8. \href{https://arxiv.org/html/2606.15624}{Primary text}.
\item
  R. Bhatia, \emph{Pinching, trimming, truncating, and averaging of
  matrices}, American Mathematical Monthly 107(7) (2000), 602--608,
  equation (2), for the complex unitary counterpart cited in reference
  1. \href{https://doi.org/10.2307/2589115}{DOI}.
\end{enumerate}

\subsection{Status check --- 2026-09-10}\label{status-check-2026-09-10}

The latest source is v2 and explicitly asks for the smallest length.
Searches included
\textit{orthogonal pinching minimum phi Bourin Lee},
\textit{Averages over matrix unitary orbits spectral order 2026},
and \textit{diagonal pinching orthogonal matrices minimum}. No
general determination was located. The adjacent Question 4.9 reverses
the quantifier order by allowing the matrices to depend on \(X\); it is
not counted separately here. This is an exact optimization question
rather than a conjectured closed formula.

\textbf{Audit update (2026-09-10):} Rechecked Question 4.8 in the
current Bourin--Lee text and searched for minimal orthogonal pinching
averages. The question remains explicit; allowing the conjugations to
depend on \(X\) is the different Question 4.9. This is a bounded
literature check, not a proof that no solution exists.
\end{document}
